\documentclass[11pt,twoside]{book}



\usepackage[T1]{fontenc}
\usepackage[utf8]{inputenc}
\usepackage{lmodern}
\usepackage[a4paper,margin=1in]{geometry}
\usepackage{amsmath, amssymb, amsthm, mathtools}
\usepackage{stmaryrd}
\usepackage{mathrsfs}
\usepackage{graphicx}
\usepackage{subcaption}
\usepackage{tikz}
\usetikzlibrary{arrows.meta, positioning, calc}
\usepackage[hidelinks]{hyperref}
\usepackage{microtype}
\usepackage{tcolorbox}
\usepackage{xcolor}
\usepackage{framed}
\usepackage[backend=biber,style=ieee]{biblatex}
\addbibresource{references.bib}
\usepackage{fontawesome5}
\usepackage{enumitem}
\usepackage{faktor}

\usepackage{thmtools}



% \declaretheoremstyle[
%   headfont=\bfseries,
%   bodyfont=\itshape,
%   spaceabove=6pt,
%   spacebelow=6pt
% ]{theoremstyle}

% \declaretheorem[style=theoremstyle, numberwithin=chapter, name=Theorem]{theorem}
% \declaretheorem[style=theoremstyle, sibling=theorem, name=Proposition]{proposition}
% \declaretheorem[sibling=theorem, name=Lemma]{lemma}
% \declaretheorem[sibling=theorem, name=Korollar]{corollary}

% \declaretheoremstyle[
%   spaceabove = 0.8em,
%   spacebelow = 0.8em,
%   headfont = \bfseries,
%   notefont = \itshape,
%   headpunct = {},
%   postheadspace = 1em,
%   headformat = \NAME~\NUMBER\NOTE\linebreak,
%   postheadspace = 10em
% ]{definitionstyle}

% \declaretheorem[style=definitionstyle, sibling=theorem, name=Definition]{definition}
% \declaretheorem[style=definitionstyle, sibling=theorem, name=Coffee]{example}

% \declaretheoremstyle[
%   headfont=\itshape,
%   bodyfont=\normalfont
% ]{remarkstyle}

% \declaretheorem[style=remarkstyle, sibling=theorem, name=Bemerkung]{remark}

% \definecolor{classical}{RGB}{153,5,5}
% \definecolor{classical_light}{RGB}{200,80,80} 



\theoremstyle{definition}
\newtheorem{Theorem}{Theorem}[chapter]
% das erste {Theorem} ist der wrapper für unten
% das zweite {Theorem} ist der Anzeigetext
% [chapter] definiert einen neuen counter der 1.1, 1.2, 2.1 etc nummeriert
\newtheorem{Definition}[Theorem]{Definition}
% {Definition} ist der wrapper für unten
% [Theorem] sagt er soll denselben counter wie für Theoreme verwenden, d.h. man hat dann Def 1.1, Def 1.2, Theorem 1.3 etc.
% {Definition} ist der Anzeigetext
\newtheorem{Example}[Theorem]{Beispiel}
\newtheorem{Corollary}[Theorem]{Korollar}
\newtheorem{Remark}[Theorem]{Bemerkung}
\newtheorem{Proposition}[Theorem]{Proposition}
\newtheorem{Lemma}[Theorem]{Lemma}
\newtheorem{Exercise}{Aufgabe}[chapter]
% das definiert einen neuen counter für Aufgaben, d.h. Def 1.1, Def 1.2, Theorem 1.3, Aufgabe 1.1
\newtheorem{Axiom}[Theorem]{Axiom}

\definecolor{shadecolor}{RGB}{238,233,233}

% der Teil kreiert jetzt farbige Boxen in die die Theoreme etc reinge-wrapped werden
\newenvironment{theorem}[1][]{\colorlet{shadecolor}{yellow!15} \begin{shaded} \begin{Theorem}[#1]}{\end{Theorem} \end{shaded}}
% \newenvironment{theorem} bestimmt wie man es später aufruft: \begin{theorem} ... \end{theorem}
% \{yellow!15} ist die Farbe der Box, bedeutet 15% yellow 85% white
% 
% [#1] ist ein Argument für das Theorem, d.h. man hat später
% \begin{theorem}[Satz von Kaffee] ... \end{theorem}
% und dann wird:
% Theorem 1.1. (Satz von Kaffee). ...
% angezeigt.
% Das Argument ist leider nicht optional, d.h. man muss jedes Mal \begin{theorem}[hier] irgendwas reinschreiben.
% \begin{theorem} oder \begin{theorem}[] kompiliert nicht.
% Ich weiß noch nicht, wie ich das optional machen kann, ohne zwei verschiedene environments zu coden.

\newenvironment{corollary}[1][]{\colorlet{shadecolor}{gray!15} \begin{shaded} \begin{Corollary}[#1]}{\end{Corollary} \end{shaded}}
\newenvironment{proposition}[1][]{\colorlet{shadecolor}{gray!15} \begin{shaded} \begin{Proposition}[#1]}{\end{Proposition} \end{shaded}}
\newenvironment{lemma}[1][]{\colorlet{shadecolor}{gray!15} \begin{shaded} \begin{Lemma}[#1]}{\end{Lemma} \end{shaded}}

\newenvironment{definition}[1][]{\colorlet{shadecolor}{blue!8} \begin{shaded} \begin{Definition}[#1]}{\end{Definition} \end{shaded}}
\newenvironment{axiom}[1][]{\colorlet{shadecolor}{blue!8} \begin{shaded} \begin{Axiom}[#1]}{\end{Axiom} \end{shaded}}

\newenvironment{remark}[1][]{\colorlet{shadecolor}{gray!15} \begin{shaded} \begin{Remark}}{\end{Remark} \end{shaded}}
\newenvironment{example}[1][]{\par \vspace{\medskipamount} \begin{Example}}{\end{Example} \par \vspace{\medskipamount}}

\newenvironment{shadedbox}{\colorlet{shadecolor}{gray!15}  \begin{shaded}{} }{ \end{shaded} }
\newenvironment{redbox}{\colorlet{shadecolor}{red!15}  \begin{shaded}{} }{ \end{shaded} }
%\newenvironment{QUESTION}{\colorlet{shadecolor}{gray!15} \begin{shaded}{Question: } }{ \end{shaded} }
% brauche ich gerade nicht
\newenvironment{exercise}{\colorlet{shadecolor}{green!12} \begin{shaded} \begin{Exercise} }{\end{Exercise} \end{shaded} }

\setlength{\parindent}{0pt}
\setlength{\parskip}{0.5em}

\hbadness=10000



\setcounter{chapter}{-1}

\renewcommand{\chaptername}{Coffee}



\newcommand{\N}{\mathbb{N}}
\newcommand{\Z}{\mathbb{Z}}
\newcommand{\Q}{\mathbb{Q}}
\newcommand{\R}{\mathbb{R}}
\newcommand{\C}{\mathbb{C}}
\newcommand{\e}{\emptyset}
\newcommand{\arr}{\longrightarrow}
\newcommand{\m}{^{-1}}
\newcommand{\Pot}{\mathcal{P}\hspace{-1pt}ot\hspace{1pt}}
\newcommand{\T}{\mathcal{T}}

\newcommand*{\defeq}{\mathrel{\vcenter{\baselineskip0.5ex \lineskiplimit0pt \hbox{\scriptsize.}\hbox{\scriptsize.}}} =}

\newcommand{\absurdum}{{\quad \huge $\lightning$}}
\renewcommand{\qedsymbol}{$\blacksquare$}
\newcommand{\subqed}{\hfill$\diamond$}


\DeclareMathOperator{\Tr}{Tr}
\DeclareMathOperator{\Det}{Det}
\DeclareMathOperator{\im}{im}
\DeclareMathOperator{\rank}{rank}



\newcounter{query}
\newenvironment{query} {\refstepcounter{query} \par\medskip \begin{redbox} \noindent \textbf{\textcolor{purple}{\thequery. }} \begingroup \color{purple} \large \ignorespaces} {\par \endgroup \end{redbox} \medskip \ignorespacesafterend}
% Das environment habe ich gecoded, um während des Schreibens Fragen, die ich nicht spontan klären oder entscheiden kann, reinzuschreiben, damit ich aufhör drüber nachzudenken und sie später aufgreifen kann. wird mit \begin{query} ... \end{query} aufgerufen. Das ist mega markant und fällt beim Korrekturlesen sofort auf.


\setlist[enumerate,1]{label=(\roman*)}


\begin{document}



\frontmatter
\tableofcontents

\mainmatter


%\chapter{Einleitung}

\citation{GURAU}

\begin{definition}[def]
    A set is a collection of objects.
\end{definition}

\begin{theorem}[trivial]
Every finite subset of $\R$ is bounded.
\end{theorem}

\begin{proof}
Basic analysis. $\sim$GURAU
\end{proof}





%\chapter{Logik}

Mathematik ist eine Geisteswissenschaft. Anders als in Naturwissenschaften wird keine Heuristik verwendet, d.h. man akzeptiert keine Aussagen, nur weil sie plausibel erscheinen oder erfahrungsgemäß stimmen. Stattdessen müssen alle Aussagen mit den Gesetzen der \emph{Logik} bewiesen werden. Der Grund dafür ist, dass die menschliche Intuition im Bereich der Mathematik nicht sehr gut ist, und man leicht etwas für richtig halten kann, das in Wirklichkeit falsch ist. Außerdem lehrt die Notwendigkeit, jede Behauptung einwandfrei beweisen zu müssen, bevor sie akzeptiert wird, eine völlig neue Denkweise, die wiederum auch in der Physik Früchte tragen kann.

Die Logik befasst sich mit dem Wahrheitsgehalt von Aussagen. In der klassischen Logik (die in der Mathematik verwendet wird) muss eine Aussage entweder wahr oder falsch sein.

\begin{example}[Aussagen]
    Beispiele für Aussagen sind:
    \begin{align*}
        & 2 + 2 = 4 & \emph{ist wahr.}\\
        & 0 = 7 & \text{\emph{ist falsch.}}\\
        & \text{Alle Menschen studieren Physik.} & \text{\emph{ist falsch.}}\\
        & \text{Kaffee ist objektiv das beste und leckerste Getränk das es gibt.} & \text{\emph{ist wahr.}}\\
    \end{align*}
\end{example}

\begin{definition}
    Sei $A$ eine Aussage. Dann ist die \emph{negierte Aussage} $\neg A$ ("nicht $A$") genau dann wahr, wenn $A$ falsch ist.
\end{definition}

\begin{definition}
    Seien $A$, $B$ Aussagen.
    Die Aussage $A \land B$ ("$A$ und $B$") ist genau dann wahr, wenn $A$ und $B$ beide wahr sind.
    Die Aussage $A \lor B$ ("$A$ oder $B$") ist genau dann wahr, wenn mindestens eine der beiden Aussagen $A$ und $B$ wahr ist.
\end{definition}

Mehrere Aussagen können durch die Verknüpfungen $\land$ (\emph{und}) und $\lor$ (\emph{oder}) aneinandergereit werden. Dabei ist  $\lor$ ist inklusiv, d.h. wenn $A \land B$ wahr ist, dann ist auch $A \lor B$ wahr. $A \lor \neg A$ ist immer wahr, und $A \land \neg A$ ist immer falsch.

Das Kernstück der Logik sind \emph{Implikationen}:

\begin{definition}
    Seien $A$, $B$ Aussagen.
    Die Aussage $A \implies B$ ("$A$ impliziert B") ist genau dann wahr, wenn daraus, dass $A$ wahr ist, folgt, dass $B$ auch wahr ist.

    $A$ nennt man Prämisse, oder hinreichende Bedingung für $B$. $B$ nennt man Schlussfolgerung, oder notwendige Bedingung für $A$.

    Die Aussage $A \iff B$ ("$A$ ist äquivalent zu $B$", "genau dann $A$ wenn $B$") ist genau dann wahr, wenn $A \implies B$ und $B \implies A$ wahr sind.
\end{definition}

Man bemerke, wenn $A$ falsch ist, dann ist $A \implies B$ wahr. Das ist nötig, weil man aus $A \implies B$ noch keinen Schluss über die Rückrichtung $B \implies A$ treffen kann, das macht erst die Äquivalenz.

\begin{example}
    Beispiele:
    \begin{align*}
        & 1 + 1 = 2 \implies 2 + 2 = 4 & \text{\emph{ist wahr, da $2 +2 = 4$ wahr ist.}}\\
        & 2 = 7 \implies 0 = 0 & \text{\emph{ist wahr, da $2 = 7$ falsch ist.}}\\
        \intertext{Man bemerke, dass man aus einer falschen Aussage alles folgern kann (hier z.B. durch Multiplikation mit $0$), das aber keinen mathematischen Wert hat: Die Implikation $2 = 7 \implies 0 = 1$ ist genauso wahr.}
        & \text{Es regnet} \implies \text{die Straße ist nass.} & \emph{ist wahr:}\\
        \intertext{Wenn es regnet, dann ist die Straße nass, also sind beide Seiten und damit auch die Implikation wahr. Wenn es nicht regnet, dann ist "es regnet" falsch, und damit die Implikation wahr.}
        & \text{Physik ist ein Studienfach} & \\
        & \implies \text{alle Menschen studieren Physik.} & \text{\emph{ist falsch:}}\\
        \intertext{"Physik ist ein Studienfach" ist wahr, aber "alle Menschen studieren Physik" ist falsch. Damit ist die Implikation falsch.}
        & 1 + 1 = 2 \iff 2 + 2 = 4 & \text{\emph{ist wahr, da beide Aussagen wahr sind.}}\\
        & 2 = 7 \iff 0 = 0 & \text{\emph{ist falsch, da $0 = 0$ wahr aber $2 = 7$ falsch ist.}}\\
        & 2 = 7 \iff 0 = 1 & \text{\emph{ist wahr, da beide Aussagen falsch sind.}}\\
    \end{align*}
\end{example}

In der Mathematik startet man mit \emph{Axiomen}, das sind Aussagen, die man als wahr annimmt oder akzeptiert. In der modernen Mathematik\footnote{Die Gödelschen Unvollständigkeitssätze besagen in etwa, dass jedes widerspruchsfreie Axiomsystem unvollständig ist, d.h. dass es immer Aussagen gibt, die man mit den gegebenen Axiomen weder beweisen noch widerlegen kann. (Das ist eine grobe Erklärung der Sätze, nicht die formale Aussage.)} sind das die Zermelo-Fränkel Axiome \ref{ZF-axioms} mit (oder selten ohne) Auswahlaxiom \ref{axiom-of-choice}. Ausgehend davon leitet man durch logisch gültige Schlüsse weitere wahre Aussagen her.

Kleine "Hilfsaussagen" ohne eigenständige Bedeutung nennt man \emph{Lemmata} (Singular \emph{Lemma}), wichtige Zwischenergebnisse nennt man \emph{Propositionen}, und ganz große Resultate nennt man \emph{Sätze} oder \emph{Theoreme}, die Grenzen sind allerdings etwas schwammig. Übliche Beweisstrategien sind:

\begin{itemize}
    \item \emph{Direkter Beweis}: Man zeigt $A \implies B$ durch eine Kette von Implikationen.
    \item \emph{Beweis durch Kontraposition}: Man zeigt $\neg B \implies \neg A$, was logisch äquivalent zu $A \implies B$ ist.
    \item \emph{Beweis durch Widerspruch}: Man nimmt an, dass $A$ wahr und $B$ falsch ist, und zeigt, dass daraus eine falsche Aussage folgt. Dabei nutzt man aus, dass $A \land \neg B$ die logische Negation von $A \implies B$ ist, und aus einer wahren Aussage keine falsche Aussage folgen kann.
\end{itemize}

\begin{definition}
    Für Aussagen $A(x)$, die von einer Variable $x$ abhängen, verwendet man \emph{Quantoren}: $\exists$ ("es existiert"), $\exists!$ ("es existiert genau ein"), und $\forall$ ("für alle"). Ein Doppelpunkt $:$ bedeutet einfach "mit ..." oder "so, dass ... gilt".
    \begin{itemize}
        \item $\exists x: A(x)$ ist wahr, wenn mindestens ein $x$ existiert, für das $A(x)$ wahr ist.
        \item $\exists! x: A(x)$ ist wahr, wenn genau ein $x$ existiert, für das $A(x)$ wahr ist.
        \item $\forall x: A(x)$ ist wahr, wenn $A(x)$ wahr für alle x ist.
    \end{itemize}
\end{definition}

\begin{remark}
    Die logischen Negationen von Quantoren sind:
    \begin{itemize}
        \item $\neg (\exists x: A(x)) \iff \forall x: \neg A(x)$
        \item $\neg (\forall x: A(x)) \iff \exists x: \neg A(x)$
    \end{itemize}
\end{remark}

\begin{remark}
    Die Beweisstrategien für Quantoren sind:
    \begin{itemize}
        \item $\exists x: A(x)$ beweist man, indem man ein $x$ explizit angibt, für das $A(x)$ wahr ist, oder die Existenz dieses $x$ aus einer wahre Existenzaussage herleitet.
        \item $\exists! x: A(x)$ beweist man in zwei Schritten. Erst zeigt man die Existenzaussage $\exists x: A(x)$. Für die Eindeutigkeitsaussage nimmt man an, dass zwei $x, y$ existieren, für die $A(x)$ und $A(y)$ wahr sind, und folgert, dass dann $x = y$ gelten muss.
        \item $\forall x: A(x)$ beweist man, indem man ein $x$ als beliebig gegeben annimmt, und dann zeigt, dass $A(x)$ (unabhängig von der Wahl von $x$) wahr ist. Möglicherweise sind hierfür Fallunterscheidungen nötig (wie z.B. Fall 1 $x = 0$ und Fall 2 $x \neq 0$ wenn man in einem Beweisschritt durch $x$ teilen will), dann muss man $A(x)$ für alle Fälle beweisen.
        \item \emph{Natürliche Induktion} ist eine nützliche Beweismethode für Aussagen $A(n)$, die von einer natürlichen Zahl $n$ abhängen. Man zeigt zuerst, dass $A(n)$ für wahlweise $n=0$ oder $n=1$ wahr ist. Dann zeigt man im zweiten Schritt, dass wenn $A(n)$ für eine beliebige natürliche Zahl $n$ wahr ist, dann auch $A(n+1)$ wahr ist. Aus beiden Schritten zusammen folgt, dass $A(n)$ für jede natürliche Zahl $n$ wahr ist.\\
        Die Beweismethode der natürlichen Induktion ist eng verknüpft mit der Definition der natürlichen Zahlen $\N$, siehe \ref{natural-numbers}. Anschaulich bedeutet beides, bei $0$ (oder $1$) anzufangen, und sukzessive mit $+1$ "hochzugehen".
    \end{itemize}
\end{remark}



\cite{GURAU}


\newpage

\fcolorbox{gray}{gray}{

    double Test.
}

%\chapter{Mengen und Zahlen}

\section{Mengen}

Mengen sind ein zentrales Konzept, auf dem die ganze Mathematik beruht. Grob gesprochen handelt es sich um eine Zusammenfassung von Objekten, die man \emph{Elemente} nennt, zu einem einzigen neuen Objekt, der \emph{Menge}. Für eine Menge $M$, die z.B. ein Element namens $x$ enthält, schreibt man $x \in M$ (sprich "x Element M" oder "x in M"). Wenn man alle Elemente einer Menge explizit kennt, beispielsweise $x$, $y$, $z$, dann kann man die Menge auch explizit durch $M = \{x, y, z\}$ mit sogenannten Mengenklammern $\{$...$\}$ oder $\{$ .. $\}$ angeben. Die Elemente können beliebige viele oder beliebig wenige sein, und sie müssen nicht intuitiv sinnvoll zusammenpassen: zum Beispiel ist $\{1, \sqrt{\pi}, \text{Gurau}, \text{\faCoffee} \}$ eine völlig valide Menge.

Mengen sind primitiv, in dem Sinne, dass sie nicht aus anderen Objekten definiert werden. Stattdessen werden sie indirekt definiert, indem mit den Gesetzen der Logik alle Eigenschaften axiomatisch aufgelistet werden, die Mengen erfüllen sollen. Das resultierende System nennt man \underline{ZF-Mengenlehre}, nach den Mathematikern Ernst Zermelo und Abraham Fraenkel. Aber zuerst definieren wir etwas Notation:\\

\begin{definition}[Teilmenge, Vereinigungsmenge, Schnittmenge]
    \label{subset-union-cap}
    \begin{enumerate}
        \item Sei $A$ eine Menge. Eine \underline{Teilmenge} $B$ von $A$ ist eine Menge, die elementweise komplett in $A$ enthalten ist, logisch definiert über
        $$x \in B \implies x \in A.$$
        Man schreibt $B \subset A$ oder $B \subseteq A$.\\
        Für eine \underline{echte Teilmenge}, d.h. $B \subset A$ und $B \neq A$, schreibt man $B \subsetneq A$.
        \item Seien $A$, $B$ Mengen. Die \underline{Vereinigung} $A \cup B$ ist:
        $$A \cup B = \bigcup \{A, B\}.$$
        (siehe \ref{ZF-axioms}.) Ihre Existenz ist durch (ZF2) und (ZF4) begründet, und damit gilt $A \cup B = \{x | x \in A \lor x \in B\}$.
        \item Der \underline{Durchschnitt} $A \cap B$ ist:
        $$A \cap B = \{x \in A | x \in B\}.$$
        Seine Existenz ist begründet durch (ZF3) mit dem Prädikat $P(x) = x \in B$, und damit gilt $A \cap B = \{x | x \in A \land x \in B\}$.
    \end{enumerate}
\end{definition}

Natürlich ergeben diese Definitionen streng genommen erst Sinn, sobald man die dafür notwendigen ZF-Axiome (ZF 1-4) postuliert hat. Wir verwenden diese Notation allerdings für eine leichter verständliche Version der übrigen ZF-Axiome (ZF 6,7,9), und haben sie daher im Sinne der Übersichtlichkeit vorangeschoben. Der Leser kann sich selbst überzeugen, dass hier kein Zirkelschluss vorliegt.

\begin{query}
    Klingt das zu gehoben? Kann ich diesen Stil im Laufe des Scriptes immer wieder rausziehen?
\end{query}

\begin{axiom}[Zermelo-Fraenkel-Axiome] %[vgl Axiom 2.1]
    \label{ZF-axioms}
    In diesem Axiomsystem wird nicht zwischen Mengen und Elementen unterschieden, d.h. jedes Element einer anderen Menge wird auch selbst als eigenständige Menge gesehen. Die einzigen Objekte, deren Existenz axiomatisiert wird, sind Mengen.
    \begin{enumerate}[label=(ZF\arabic*)]
        \item \underline{Extensionalitätsaxiom}: Eine Menge wird vollständig durch ihre Elemente bestimmt. Formal, für je zwei Mengen $A$, $B$ gilt $(x \in A \iff x \in B) \implies A = B$. Das Extensionalitätsaxiom impliziert unter anderem, dass weitere im folgenden axiomatisierten Mengen eindeutig sind.
        \item \underline{Paarmengenaxiom}: Wenn $A$, $B$ Mengen sind, dann ist $\{A, B\}$ eine Menge.
        \item \underline{Aussonderungsaxiom}: Wenn $A$ eine Menge ist, und $P(x)$ eine Eigenschaft\footnote{Logisch streng genommen ist $P$ ein \emph{Prädikat}.}, die Elemente $x$ von $A$ entweder erfüllen oder nicht erfüllen können, dann ist $\{x \in A | P(x)\}$ eine Menge. Die Notation $x \in A | P(x)$ bedeutet "x in A, für die gilt: P(x)". Logisch ist das wieder definiert durch $x \in A \land P(x)$, aber der senkrechte Strich trennt optisch die Elemente der Menge von den Bedingungen, die diese erfüllen.
        \item \underline{Vereinigungsaxiom}: Wenn $A$ eine Menge ist, die aus Mengen besteht, dann ist $\bigcup A \defeq \{x | x \in B \land D \in A\}$ eine Menge.
        \item \underline{Leermengenaxiom}: Es existiert eine Menge, die keine Elemente enthält. Diese nennen wir die \underline{leere Menge} und schreibt $\e$ oder $\{\}$. Formal ist die Aussage $\exists \e: \forall x: x \notin \e$. Die leere Menge ist die einzige explizite Menge, deren Existenz axiomatisiert wird. Alle weiteren expliziten Mengen müssen also mit Hilfe der anderen Axiome aus der leeren Menge konstruiert werden.
        \item \underline{Potenzmengenaxiom}: Für jede Menge $A$ existiert eine Menge, die alle Teilmengen von $A$ enthält. Diese nennen wir die Potenzmenge $\mathscr{P}(A)$. Dann ist $\mathscr{P}(A) = \{B | B \subset A\}$.
        \item \underline{Unendlichkeitsaxiom}: Sei $A$ eine beliebige Menge. Dann existiert eine Menge $B$ mit $\e \in B$, und $A \in B \implies A \cup \{A\} \in B$.\\
        Man bemerke, dass mit $A \arr A \cup \{A\} \in B$ die Voraussetzung für das Unendlichkeitsaxiom wieder erfüllt ist, damit gilt also auch $A \cup \{A\} \cup \{A \cup \{A\}\} \in B$. Diese Kette lässt sich \underline{induktiv}, d.h. durch wiederholtes Ausführen desselben Argumentes, beliebig weit fortsetzen. $B$ enthält also, sofern es mindestens ein $A$ enthält, gleich unendlich viele Mengen.
        \item \underline{Ersetzungsaxiom}: Ist $A$ eine Menge, und $f$ eine Vorschrift, die jede Menge $x \in A$ eindeutig durch eine andere Menge $f(x)$ ersetzt (d.h. $f$ ist eine \emph{Funktion}, Def. \ref{def-function}), dann ist das Bild von $f$ definiert als $f(A) = \{f(x) | x \in A\}$ wieder eine Menge.
        \item \underline{Fundierungs}- oder \underline{Regularitätsaxiom}: Sei $A$ eine Menge. Dann enthält $A$ ein Element $x$ mit $x \cap A = \e$.
    \end{enumerate}
\end{axiom}

Danach ein paar Bemerkungen zu einzelnen Axiomen

\begin{remark}
    \begin{itemize}
        \item e
    \end{itemize}
\end{remark}

Danach noch ein paar Definitionen

\begin{definition}[Komplement, symmetrische Differenz, kartesisches Produkt]
    Seien $A$, $B$ Mengen. Wir definieren
    \begin{enumerate}
        \item das \underline{Komplement} $B \backslash A = \{x \in B | x \notin A\}$,
        \item die \underline{symmetrische Differenz} $A \triangle B = (A \backslash B) \cup (B \backslash A)$,
        \item das \underline{kartesische Produkt} $A \times B = \{ (a, b) | a \in A, b \in B\}$.
    \end{enumerate}
\end{definition}

Man bemerke, dass die Paare $(a, b)$ geordnet sind, d.h. wenn $A$, $B$ zwei Elemente $x \neq y$ gemeinsam haben, dann ist $(x, y) \neq (y, x)$. Die Existenz von geordneten Paaren ist z.B. durch (ZF2) mittels einer Identifikation wie $(a, b) = \{\{a\}, \{a, b\}\}$ begründet.

\begin{example}
    zu Mengenoperationen $\subset, \cup, \cap, \backslash, \triangle, \times$ auch Beispiele mit $\e$!
\end{example}

\begin{lemma}[Regeln von De Morgan]
    Seien $A$, $B$, $C$ Mengen. Dann gilt:
    \begin{enumerate}
        \item $A \backslash (B \cup C) = (A \backslash B) \cap (A \backslash C)$,
        \item $A \backslash (B \cap C) = (A \backslash B) \cup (A \backslash C)$.
    \end{enumerate}
\end{lemma}
Wir sehen nun unseren ersten kleineren Beweis. Die Beweismethode, die wir hier verwenden, nennt man \emph{Nachrechnen}, damit meint man in der Mathematik das wiederholte Anwenden von Axiomen, Definitionen, und bereits gezeigten Aussagen, ohne dabei größere Ideen zu brauchen oder auf kompliziertere Details achten zu müssen.
\begin{proof}
    \begin{enumerate}
        \item Sei $x$ beliebig. Man bemerke, dass $x \notin A \iff \neg (x \in A)$ ist. Dann gilt mit den Gesetzen der Logik:\\
        $x \in A \backslash (B \cup C)$\\
        $\iff x \in A \land x \notin (B \cup C)$\\
        $\iff x \in A \land (x \notin B \land x \notin C)$\\
        $\iff (x \in A \land x \notin B) \land (x \in A \land x \notin C)$\\
        $\iff x \in (A \backslash B) \cap (A \backslash C)$.\\
        Nach (ZF1) sind also beide Mengen gleich.
        \item Analog zu (i). Sei $x$ beliebig. Dann gilt:\\
        $x \in A \backslash (B \cap C)$\\
        $\iff x \in A \land x \notin (B \cap C)$\\
        $\iff x \in A \land (x \notin B \lor x \notin C)$\\
        $\iff (x \in A \land x \notin B) \lor (x \in A \land x \notin C)$\\
        $\iff x \in (A \backslash B) \cup (A \backslash C)$.\\
        Nach (ZF1) sind also beide Mengen gleich.
    \end{enumerate}
\end{proof}




\section{Relationen}

\begin{definition}[Relation]
    Seien $A$, $B$ Mengen. Eine \underline{Relation} $\sim_r$ zwischen $A$ und $B$ ist eine Zuordnungsvorschrift, die zwei Elemente $a \in A$ und $b \in B$ eindeutig entweder in Relation zueinander ($a \sim_r b$ ist wahr), oder nicht in Relation zueinander ($a \sim_r b$ ist falsch) stellt.\\
    Man identifiziert dabei die Relation als Teilmenge $R \subset A \times B$, indem man
    $$a \sim_r b \text{ ist wahr, oder einfach } a \sim_r b \iff (a,b) \in R \subset A \times B$$
    setzt.\\
    Wenn $A = B$ ist, dann nennt man $\sim_r$ auch einfach eine Relation auf $A$.
\end{definition}

Ihre Existenz erhalten Relationen dadurch, dass sie formal Teilmengen sind. Durch Forderung von weiteren Eigenschaften haben Relationen unterschiedlichste Anwendungen.

Auf Wikipedia unter \href{https://de.wikipedia.org/wiki/Relation_(Mathematik)#Beispiele}{Relationen} (und verlinkte Artikel) findet man sehr anschauliche Beispiele für Relationen (größtenteils mit fragwürdiger mathematischer Relevanz). Es gibt aber auch viele anwendungsbezogene Beispiele für Relationen, eines der wichtigsten davon ist die Äquivalenzrelation.

\begin{definition}[Äquivalenzrelation] %vgl Def 2.1
    Sei $X$ eine Menge. Eine Relation $\sim$ auf $X$ heißt \underline{Äquivalenzrelation}, falls folgende Eigenschaften gelten:
    \begin{enumerate}[label=(Ä\arabic*)]
        \item \underline{reflexiv}: $\forall x \in X: x \sim x$
        \item \underline{symmetrisch}: $\forall x, y \in X: x \sim y \iff y \sim x$
        \item \underline{transitiv}: $\forall x, y, z \in X: x \sim y \land y \sim z \implies x \sim z$
    \end{enumerate}
\end{definition}

\begin{definition}[Äquivalenzklassen]
    \label{def-equivalence-class}
    Sei $X$ eine Menge und $\sim$ eine Äquivalenzrelation auf $X$. Für ein festes $x \in X$ definieren wir die \underline{Äquivalenzklasse} von $x$
    $$[x] = \{y | x \sim y \},$$
    und die \underline{Quotientenmenge} von $\sim$
    $$\faktor{X}{\sim} = \{ [x] | x \in X \}.$$
\end{definition}

\begin{lemma}[Äquivalenzklassen sind wohldefiniert]
    Sei $X$ eine Menge mit Äquivalenzrelation $\sim$.
    
    Dann sind die Äquivalenzklassen, d.h. die Elemente von $\faktor{X}{\sim}$, nichtleer, abgeschlossen, und disjunkt.
\end{lemma}
Wir führen diesen Beweis sehr ausführlich durch, um zu zeigen, wie man mathematisch rigoros arbeitet.
\begin{proof}
    Sei $x \in X$.\\

    \underline{nichtleer}: Da $x \in X$, gilt nach (Ä1) $x \sim x$ und damit nach Def. \ref{def-equivalence-class} $x \in [x]$.
    \pushQED{\subqed}\popQED\\

    Sei nun auch $y \in X$.

    \underline{abgeschlossen}: Das bedeutet, wenn $y \in [x]$, dann ist $[x] = [y]$, d.h. jede Äquivalenzklasse bleibt "unter sich".

    Sei also $y \in [x]$.\\
    Wir zeigen beide Mengeninklusionen $[x] \subset [y]$ und $[y] \subset [x]$, daraus folgt $[x] = [y]$ nach (ZF1).

    "$\subset$": Die Annahme ist $y \in [x]$.\\
    Nach der Definition von Äquivalenzklassen \ref{def-equivalence-class} bedeutet das $x \sim y$.\\
    Wegen Symmetrie (Ä2) folgt dann auch $y \sim x$.

    Sei nun weiter $z \in [x]$.\\
    Nach \ref{def-equivalence-class} bedeutet das wieder $x \sim z$.
    
    Da wir nun $y \sim x$ und $x \sim z$ haben, gilt mit Transitivität (Ä3) $y \sim z$.\\
    Das ist aber nach \ref{def-equivalence-class} wieder die Definition von $z \in [y]$.
    
    Insgesamt haben wir $y \in [x] \implies y \in [z]$ gezeigt.\\
    Es gilt also nach \ref{subset-union-cap} $[x] \subset [y]$.

    Sei nun $z \in [y] \implies y \sim z$, da $x \sim y \implies x \sim z \implies z \in [x]$.
    
    Es gilt also auch $[y] \subset [x]$, und insgesamt damit $[x] = [y]$.
    \pushQED{\subqed}\popQED

    \underline{disjunkt}: Sei wieder $y \in X$, dann ist $y \in [y]$.

    Zu zeigen ist, dass wenn $x \nsim y$ (d.h. $y \notin [x]$), dann folgt $[x] \cap [y] = \e$.

    Wir zeigen die Kontraposition: $[x] \cap [y] \neq \e \implies x \sim y$.

    Sei nun $[x] \cap [y] \neq \e$, dann $\exists z \in [x] \cap [y]$. Wähle so ein $z \in [x] \cap [y]$ beliebig.

    Mit $z \in [x] \implies x \sim z$, und mit $z \in [y] \implies y \sim z \implies z \sim y \implies x \sim z$.
\end{proof}

\begin{query}
    Wie sind die Zeilenabstände in dem Beweis? Ist leicht lesbar so, aber das Dokument wird schon ziemlich lang ($\sim$ paar hundert Seiten nur für Hoema 1) werden wenn ich das durchziehe.

    Ein alternativer Vorschlag, der die Länge nicht wirklich verändert, aber dafür zusammengehörende Zeilen enger rückt auf Kosten der Lesbarkeit:
\end{query}

\begin{example}
    Äquivalenzklassen
    \begin{enumerate}
        \item Restklassen
    \end{enumerate}
\end{example}

\begin{definition}[Halbordnungsrelation, halbgeordnete Menge, Totalordnungsrelation]
    Sei $X$ eine Menge. Eine Relation $\preceq$ auf $X$ heißt \underline{Halbordnungsrelation} (oder \underline{Ordnungsrelation}), wenn gilt:
    \begin{enumerate}[label=(H\arabic*)]
        \item \underline{reflexiv}: $\forall x \in X: x \preceq x$
        \item \underline{antisymmetrisch}: $\forall x, y \in X: x \preceq y \land y \preceq x \implies x = y$
        \item \underline{transitiv}: $\forall x, y, z \in X: x \preceq y \land y \preceq z \implies x \preceq z$
    \end{enumerate}
    Dann heißt das Tupel $(X, \preceq)$ \underline{halbgeordnete Menge} (\underline{Poset} von engl. "partially ordered set")\footnote{Endliche Posets können durch sogenannte Hasse-Diagramme dargestellt werden, das sind flow charts, in denen Elemente durch Knotenpunkte, und Relationen $x \preceq y$ durch Pfeile von $x$ nach $y$ dargestellt werden.}.

    Eine Halbordnungsrelation heißt \underline{Totalordnungsrelation}, falls zusätzlich gilt:
    \begin{enumerate}[label=(T)]
        \item $\forall x, y \in X: x \preceq y \lor y \preceq x$
    \end{enumerate}
    Eine nichtleere, total geordnete Teilmenge einer halbgeordneten Menge heißt \underline{Kette}.
\end{definition}

\begin{example}
    \label{exa-posets}
    \begin{itemize}
        \item $(\R, \leq)$ ist ein Poset, sogar total
        \item $(X, \preceq)$ mit $x \preceq y \iff x = y$ ist ein Poset, aber nicht total
        \item "ganzzahliger Teiler von"
        \item $A \subset X$ ist poset
    \end{itemize}
\end{example}

\begin{query}
    Warum trinkst Du immer noch keinen Kaffee?
\end{query}

\begin{definition}[strenge Halbordnungsrelation, strenge Totalordnungsrelation]
    Sei $X$ eine Menge. Eine Relation $\prec$ auf $X$ heißt \underline{strenge Halbordnungsrelation}, wenn gilt:
    \begin{enumerate}[label=(\arabic*)]
        \item \underline{irreflexiv}: $\forall x \in X: x \not\prec x$
        \item \underline{asymmetrisch}: $\forall x, y \in X: x \prec y \implies y \not\prec x$
        \item \underline{transitiv}: $\forall x, y, z \in X: x \prec y \land y \prec z \implies x \prec z$
    \end{enumerate}

    $\prec$ heißt \underline{strenge Totalordnungsrelation} auf $X$, wenn zusätzlich gilt:
    \begin{enumerate}[label=(\arabic*), start=4]
        \item $\forall x, y \in X$ gilt genau eine der folgenden drei Eigenschaften: $x \prec y$, oder $y \prec x$, oder $x = y$
    \end{enumerate}
\end{definition}

\begin{example}
    strenge Halbordnungen und Totalordnungen
\end{example}

\begin{definition}[obere Schranke, Maximum]
    \label{def-upper-boundary}
    Sei $(X, \preceq)$ eine halbgeordnete Menge, und $A \subset X$. Dann heißt $a \in X$
    \begin{enumerate}
        \item \underline{obere Schranke} von $A$, wenn $\forall x \in A: x \preceq a$.
        
        Gilt zusätzlich $a \in A$, dann heißt $a$ \underline{größtes Element} oder \underline{Maximum} von $A$. In diesem Fall schreiben wir $a = \max(A)$.
        \item \underline{untere Schranke} von $A$, wenn $\forall x \in A: a \preceq x$.
        
        Gilt zusätzlich $a \in A$, dann heißt $a$ \underline{kleinstes Element} von $A$. In diesem Fall schreiben wir $a = \min(A)$.
    \end{enumerate}
\end{definition}

\begin{remark}
    Im Allgemeinen sind (hier und in den folgenden Definitionen) beide Konzepte identisch. Zu jeder Halbordnung kann man nämlich eine umgedrehte Halbordnung $\succeq$ definieren, indem man $x \succeq y \iff y \preceq x$ setzt. Dann ist die erste Bedingung $a \succeq x$ und die zweite Bedingung $x \succeq a$. Daraus sieht man sofort, dass eine obere Schranke bezüglich $\preceq$ eine untere Schranke bezüglich $\succeq$ ist, etc.
\end{remark}

\begin{example}
    boundaries, Min, Max
\end{example}

\begin{lemma}[lemma]
    \label{lem-upper-bound-unique}
    Sei $(X, \preceq)$ eine halbgeordnete Menge.
    
    Falls eine obere Schranke von $X$ existiert, dann ist diese eindeutig.
\end{lemma}
\begin{proof}
    Seien $a, b \in X$ obere Schranken von $X$.

    $a$ ist obere Schranke $\implies \forall x \in X: x \preceq a$. Wenn wir $x = b$ einsetzen $\implies b \preceq a$.
    
    $b$ ist obere Schranke $\implies \forall x \in X: x \preceq b$. Wenn wir $x = a$ einsetzen $\implies a \preceq b$.

    Aus beiden Aussagen zusammen folgt mit (H2) $\implies a = b$.
\end{proof}

Analog ist auch eine untere Schranke von $X$, falls sie existiert, eindeutig.

\begin{remark}
    \label{rem-upper-bound-unique}
    Wir bemerken, dass sich diese Aussage nur auf eine Schranke der kompletten Menge bezieht. Schranken von $A \subset X$ müssen nicht eindeutig sein, da sie nicht in $A$ liegen müssen:
    
    Zum Beispiel sind in $(\R, \leq)$ mit $[0, 1] \subset \R$ sowohl $1 \in \R$ als auch $2 \in \R$ (und analog jede reelle Zahl $\geq 1$) obere Schranken von $[0, 1]$, aber es gibt nur eine obere Schranke von $[0, 1]$, die in $[0, 1]$ liegt, nämlich $1$.

    Falls es allerdings eine obere Schranke von $A$ gibt, die in $A$ liegt, dann ist diese wieder eindeutig, da $(A, \preceq)$ nach \ref{exa-posets} wieder eine halbgeordnete Menge ist.
\end{remark}

\begin{definition}[maximales Element]
    \label{def-maximal-element}
    Sei $(M, \preceq)$ eine halbgeordnete Menge, und $A \subset M$. Dann heißt $a \in A$
    \begin{enumerate}
        \item \underline{maximales Element} von $A$, falls $\forall x \in A: a \preceq x \implies a = x$.
        \item \underline{minimales Element} von $A$, falls $\forall x \in A: x \preceq a \implies a = x$.
    \end{enumerate}
\end{definition}

Ein maximales Element ist eine schwächere Bedingung als eine obere Schranke, da anders als in \ref{def-upper-boundary} keine Vergleichbarkeit gefordert wird. Stattdessen muss ein maximales Element nur größer als alle Elemente sein, mit denen es vergleichbar ist. Ein einfaches Beispiel verdeutlicht den Unterschied:

\begin{example}
    $X = \{ \e, \{1\}, \{2\}, \{3\}, \{1, 2\}, \{1, 3\}, \{2, 3\}\}$ mit der Relation $\subset$.

    Dann sind $\{1, 2\}$, $\{1, 3\}$, und $\{2, 3\}$ maximale Elemente, da sie jedes andere Element entweder enthalten, z.B. $\e \subset \{1, 2\}$ und $\{2\} \subset \{1, 2\}$, oder nicht mit ihm vergleichbar sind, z.B. $\{1, 2\} \not\subset \{1, 3\}$ und $\{1, 3\} \not\subset \{1, 2\}$.

    An diesem Beispiel sehen wir auch, dass maximale Elemente aufgrund fehlender Vergleichbarkeit nicht eindeutig sein müssen.

    Wenn wir aber noch $\{1, 2, 3\}$ zur Menge $X$ hinzufügen, dann gibt es nur noch ein maximales Element, nämlich $\{1, 2, 3\}$. Dieses ist auch eine obere Schranke, da es alle andere Elemente enthält.
    
    Dieser Zusammenhang gilt übrigens auch im Allgemeinen:
\end{example}

\begin{lemma}[lemma]
    \label{lem-bound-max-el}
    Sei $(M, \preceq)$ eine halbgeordnete Menge, $A \subset M$, und $a \in M$.

    Wenn $a$ obere Schranke von $A$ ist, dann ist $a$ maximales Element von $A$.
\end{lemma}
\begin{proof}
    Sei $a \in M$ obere Schranke von $A$.
    
    Zu zeigen ist, dass $\forall x \in A$ gilt: $a \preceq x \implies a = x$.

    Sei also $x \in A$ mit $a \preceq x$.
    
    Da $a$ obere Schranke von $A$ ist, gilt auch $x \preceq a \forall x \in A$.

    Mit Antisymmetrie (H2) folgt $a = x$, also ist $a$ maximales Element von $A$.
\end{proof}

\begin{lemma}[lemma]
    \label{lem-maximum-max-el}
    Sei $(M, \preceq)$ eine halbgeordnete Menge, $a \in M$, und $A \subset M$.

    Wenn $a$ Maximum von $A$ ist, dann ist $a$ das eindeutige maximale Element von $A$.
\end{lemma}
\begin{query}
    Das wäre eine gute Aufgabe für den Leser?
\end{query}
\begin{proof}
    Sei $a$ Maximum von $A$. Dann ist $a$ nach \ref{def-upper-boundary} auch obere Schranke von $A$. Nach \ref{lem-bound-max-el} ist $a$ maximales Element von $A$.

    Als Maximum muss $a$ wiederum per Definition in $A$ liegen, und als obere Schranke von $A$ in $A$ ist damit $a$ nach \ref{lem-upper-bound-unique} / \ref{rem-upper-bound-unique} eindeutig.
\end{proof}

\begin{example}
    untere, obere schr
\end{example}

\begin{definition}[Infimum, Supremum]
    Sei $(M, \preceq)$ eine halbgeordnete Menge, $a \in M$, und $A \subset M$.
    \begin{enumerate}
        \item Ist $a$ obere Schranke von $A$, dann heißt $a$ \underline{Supremum} von $A$, falls $\forall x \in M$ gilt: $x$ ist obere Schranke von $A \implies a \preceq x$. In diesem Fall schreiben wir $a = \sup(A)$.
        \item Ist $a$ untere Schranke von $A$, dann heißt $a$ \underline{Infimum} von $A$, falls $\forall x \in M$ gilt: $x$ ist untere Schranke von $A \implies x \preceq a$. In diesem Fall schreiben wir $a = \inf(A)$.
    \end{enumerate}
\end{definition}

\begin{example}
    inf, sup
\end{example}

\begin{definition}[induktiv geordnet]
     Eine eine halbgeordnete Menge $(M, \preceq)$ heißt \underline{induktiv geordnet}, wenn zu jeder Kette in $M$ eine obere Schranke in $M$ existiert.
\end{definition}

\begin{example}
    $Pot(X), \subset$ induktiv
    
    $\N, \leq$ nicht induktiv

    $[0, 1], [0, 1)$
\end{example}

\begin{definition}[Wohlordnungsrelation]
    Sei $X$ eine Menge mit Totalordnungsrelation $\preceq$. Dann heißt $\preceq$ \underline{Wohlordnungsrelation} und $(X, \preceq)$ \underline{wohlgeordnete Menge}, falls zusätzlich gilt:
    \begin{enumerate}[label=(W)]
        \item In jeder nichtleeren Teilmenge von $X$ existiert ein Minimum bezüglich $\preceq$.
    \end{enumerate}
\end{definition}

\begin{remark}
    Eine Wohlordnung kann man als eine \emph{kanonische} (d.h. für die menschliche Intuition die natürlichste Wahl) Reihenfolge von Elementen in $X$ verstehen: Für die Teilmenge $X \subset X$ gibt (W) einem ein Minimum $w_1$, dieses nimmt man als erstes Element in der Reihenfolge. Dann gibt (W) einem für $X \backslash \{w_1\} \subset X$ wieder ein Minimum $w_2$, dann für $X \backslash \{w_1, w_2\}$ ein Minimum $w_3$, und so weiter. Durch Induktion\footnote{Endliche, natürliche, oder transfinite Induktion, je nach Kardinalität von $X$, siehe \ref{REF}.} kann man alle Elemente von $X$ relativ zueinander ordnen: $w_1, w_2, w_3, ...$

    Die Existenz einer Wohlordnung ist für endliche Posets trivial: Man zählt die Elemente einfach in einer beliebigen Reihenfolge auf, und definiert sich entsprechend $x < y \iff x$ in der Aufzählung vor $y$ steht. Die gewählte Aufzählung ist dann auch schon die Wohlordnung. Für gewisse unendliche Mengen kann man auch eine Wohlordnung angeben, z.B. für $(\N, <)$ ist die Aufzählung $\{0, 1, 2, 3, ...)\}$ eine mögliche Wohlordnung. $(\Z, <)$ und $(\Q, <)$ können auch wohlgeordnet werden, siehe dazu den Beweis von \ref{theorem-cantor-1}. Die Existenz einer Wohlordnung für alle unendlichen Mengen benötigt ein zusätzliches Axiom, siehe \ref{chapter-tenth-axiom}.
\end{remark}



\section{Funktionen}

\begin{definition}[Funktion] %vlg Def 2.1
    \label{def-function}
    Seien $X$, $Y$ Mengen. Eine \emph{Abbildungsvorschrift} $f: X \arr Y$ ist eine Vorschrift, die Elemente $x$ in $X$ nimmt, und ihnen Elemente $f(x)$ in $Y$ zuordnet. Diese Zuordnung kann erstmal beliebig sein. Eine Abbildungsvorschrift heißt \emph{wohldefiniert}, wenn sie folgende zwei Eigenschaften erfüllt:
    \begin{itemize}
        \item \underline{linkstotal}: $\forall x \in X \exists y \in Y: f(x) = y$, das heißt jedes $x$ in $X$ wird auf ein $y$ abgebildet, das auch in $Y$ liegt.
        \item \underline{rechtseindeutig}: $\forall x \in X \forall y_1, y_2 \in Y: f(x) = y_1 \land f(x) = y_2 \implies y_1 = y_2$. Das bedeutet jedes $x$ wird nur auf ein $y$ abgebildet, und nicht auf zwei verschiedene (oder was die logische Aussage bedeutet, wenn es zwei $y$s gibt auf die $x$ abgebildet wird, dann sind beide das Gleiche $y$).
    \end{itemize}
    Eine wohldefinierte Abbildungsvorschrift $f: X \arr Y$ heißt \underline{Funktion} (oder \underline{Abbildung}).
    
    $X$ heißt \underline{Definitionsmenge} (\underline{Definitionsbereich}) von $f$, $Y$ heißt \underline{Bildmenge} (\underline{Zielmenge}) von $f$.

    Die Menge $\{(x, f(x)) | x \in X\}$ heißt \underline{Graph} von $f$.
\end{definition}

\begin{remark}
    Funktionen sind keine fundamental neuen Objekte, sondern werden aus Relationen konstruiert: Formal ist eine Funktion $f$ über eine Relation $R_f \subset X \times Y$ mittels der Identifikation $f(x) = y \iff (x, y) \in R_f$ definiert. $R_f$ ist der Graph von $f$. Damit sind Funktionen also im Grunde wieder Mengen, und benötigen keine zusätzlichen Axiome, um ihre Existenz zu begründen.
\end{remark}

\begin{example}
    \begin{enumerate}[label=(\arabic*)]
        \item Auf den zwei Mengen $X = \{1, \sqrt{\pi}, \text{Gurau}, \text{\faCoffee} \}$, $Y = \{\text{rational}, \text{irrational},$
        $\text{Vorlesung}, \text{trinken}\}$ definieren wir zwei bijektive Funktionen:
        
        $f: X \arr Y$ über
        \begin{itemize}[label=]
            \item $f(1) = \text{rational}$
            \item $f(\sqrt{\pi}) = \text{irrational}$
            \item $f(\text{Gurau}) = \text{Vorlesung}$
            \item $f(\text{\faCoffee}) = \text{trinken}$
        \end{itemize}
        und $g: X \arr Y$ über
        \begin{itemize}[label=]
            \item $g(1) = \text{Vorlesung}$
            \item $g(\sqrt{\pi}) = \text{rational}$
            \item $g(\text{Gurau}) = \text{trinken}$
            \item $g(\text{\faCoffee}) = \text{irrational}$
        \end{itemize}
        Für uns Menschen sieht die Funktion $f$ nach einer sinnvollen Zuordnung aus, die Funktion $g$ dagegen völlig unsinnig. Aus mathematischer Sicht sind allerdings beide Funktionen absolut gleichwertig, nämlich wohldefiniert und bijektiv.\\
        Oft haben mathematische Konstrukte eine Motivation, die auf nicht streng logischer Intuition beruht oder aus der weltlichen Anschauung kommt, aber das ist eben nicht immer so. Und auf keinen Fall darf man Mathematik danach beurteilen, ob sie intuitiv Sinn ergibt, sondern nur ob sie logisch Sinn ergibt.
    \end{enumerate}
\end{example}

\begin{definition}[Bild, Urbild]
    Seien $X$, $Y$ Mengen und $f: X \arr Y$ eine Funktion. Seien weiter $U \subset X$ und $V \subset Y$ Teilmengen.
    \begin{enumerate}
        \item Das \underline{Bild} von $U$ unter $f$ ist die Menge $f(U) = \{f(x) \in Y | x \in U\}$.
        \item Das \underline{Urbild} von $V$ unter $f$ ist die Menge $f\m(V) = \{x \in X | \exists y \in V: f(x) = y\}$.
    \end{enumerate}
\end{definition}

\begin{definition}[neue Funktionen]
    \begin{enumerate}
        \item Die \underline{Einschränkung} von $f$ auf $U$ ist die Funktion $f|_U: U \arr Y$ definiert durch $f|_U(x) = f(x)$, $x \in U$.
        \item kanonische Einbettung
        \item Die \underline{Identität} auf einer beliebigen Menge $X$ ist die Funktion $id_X: X \arr X$ definiert durch $id_X(x) = x \forall x \in X$.\\
        Sei weiter $Z$ eine Menge und $g: Y \arr Z$ eine Funktion.
        \item Die \underline{Komposition} von $f$ und $g$ ist die Funktion $g \circ f: X \arr Z$ definiert durch $(g \circ f)(x) = f(g(x))$, $x \in X$.
    \end{enumerate}
\end{definition}

\begin{lemma}[Diese Funktionen sind wohldef]
    trivial
\end{lemma}

\begin{itemize}
    \item Indexmenge, Familie von Elementen / Funktionen
    \item Verein, Schnitte
    \item Verhalten unter $f$
\end{itemize}

\begin{example}
    möglicherweise selbe Beispiele wie in der Bem davor nochmal aufgreifen
\end{example}

\begin{definition}[injektiv, surjektiv, bijektiv]
    Seien $X$, $Y$ Mengen, und $f: X \arr Y$ eine Funktion.
    \begin{enumerate}
        \item $f$ heißt \underline{injektiv}, wenn gilt: $\forall x, y \in X: x \neq y \implies f(x) \neq f(y)$.\\
        Injektivität bedeutet, dass die Funktion jedes Element in der Zielmenge höchstens einmal "trifft".\\
        Oft lässt sich Injektivität leichter durch die Kontraposition der Aussage zeigen: $\forall x, y \in X: f(x) = f(y) \implies x = y$.
        \item $f$ heißt \underline{surjektiv}, wenn gilt: $\forall y \in Y \exists x \in X: f(x) = y$.\\
        Surjektivität bedeutet, dass die Funktion jedes Element in der Zielmenge mindestens einmal "trifft".
        \item $f$ heißt \underline{bijektiv}, wenn $f$ injektiv und surjektiv ist.\\
        Bijektivität bedeutet also, dass die Funktion jedes Element in der Zielmenge genau einmal "trifft", und damit eine eins-zu-eins-Abbildung ist.
    \end{enumerate}
\end{definition}

\begin{example}
    je mehr Bsp desto besser. wie wäre es mit ein paar Bildern?
\end{example}

\begin{remark}
    Offensichtlich kann eine Funktion $f: X \arr Y$ zwischen zwei endlichen Mengen nur...
    \begin{itemize}
        \item injektiv sein, wenn $X$ nicht mehr Elemente als $Y$ hat.
        \item surjektiv sein, wenn $Y$ nicht mehr Elemente als $X$ hat.
        \item bijektiv sein, wenn $X$ und $Y$ genau gleich viele Elemente haben.
    \end{itemize}
    Für unendliche Mengen muss man etwas vorsichtiger sein, weil es nichttrivial ist zu sagen, "wie viele" Elemente zwei unendliche Menge im Vergleich zueinander haben. Später werden wir, motiviert durch den Fall von endlichen Mengen, die Existenz von bijektiven Funktionen verwenden, um zu definieren, dass zwei unendliche Mengen gleich groß sind.
\end{remark}

\begin{lemma}[lemma]
    Seien $f: X \arr Y$, $g: Y \arr Z$ Funktionen. Dann gilt:
    \begin{enumerate}
        \item $f$ und $g$ sind injektiv $\implies g \circ f$ ist injektiv.
        \item $f$ und $g$ sind surjektiv $\implies g \circ f$ ist surjektiv.
        \item $f$ und $g$ sind bijektiv $\implies g \circ f$ ist bijektiv.
    \end{enumerate}
\end{lemma}

\begin{definition}[Umkehrfunktion]
    Seien $X$, $Y$ Mengen, und $f: X \arr Y$ eine Funktion. Sei weiter $g: Y \arr X$ eine Funktion.
    \begin{enumerate}
        \item $g$ heißt \underline{linksinverse Funktion} von $f$, wenn $g \circ f = id_X$.
        \item $g$ heißt \underline{rechtsinverse Funktion} von $f$, wenn $f \circ g = id_Y$.
        \item $g$ heißt \underline{inverse Funktion} (oder \underline{Umkehrabbildung}) von $f$, wenn $g$ linksinverse und rechtsinverse Funktion von $f$ ist. In diesem Fall nennen wir $f$ \underline{invertierbar}.
        
        Falls die inverse Funktion existiert, bezeichnen wir diese auch $f\m$.
    \end{enumerate}
\end{definition}

Inverse müssen nicht immer existieren.

Bem: Urbild / Umkehrabb!

\begin{example}
    hier wäre Bsp für Fkt die nur in eine Richtung invertierbar wären interessant. Aber auch bijektiv und gar nicht invertierbar
\end{example}

\begin{lemma}[Existenz von Umkehrabbildungen]
    \label{lem-existence-inverse-function}
    Sei $f: X \arr Y$ eine Funktion. Dann gilt:
    \begin{enumerate}
        \item $f$ ist injektiv $\iff f$ ist linksinvertierbar.
        \item $f$ ist surjektiv $\iff f$ ist rechtsinvertierbar.
        \item $f$ ist bijektiv $\iff f$ ist invertierbar.
    \end{enumerate}
\end{lemma}

\section{Das zehnte Axiom}

\begin{axiom}[Lemma von Zorn]
    \label{lemma-von-zorn}
    In jeder nichtleeren, induktiv geordneten Menge existiert ein maximales Element.
\end{axiom}

\begin{axiom}[Wohlordnungssatz]
    \label{wohlordnungssatz}
    Auf jeder Menge existiert eine Wohlordnungsrelation.
\end{axiom}

\begin{axiom}[Auswahlaxiom]
    \label{axiom-of-choice}
    Ist $I$ eine Indexmenge, und $(X_i)_{i \in I}$ eine Familie von nichtleeren, paarweise disjunkten Mengen, dann existiert eine Abbildung $f: I \arr \cup_{i \in I}X_i$ mit $f(i) \in X_i \forall i \in I$.
\end{axiom}

\begin{theorem}[C]
    Das Lemma von Zorn, der Wohlordnungssatz, und das Auswahlaxiom sind innerhalb der ZF-Mengenlehre äquivalent.
\end{theorem}
\begin{proof}
    trivial $\sim GURAU$
\end{proof}

\begin{query}
    Wäre hier das Unterkapitel über Gruppen, Halbgruppen etc sinnvoll? Oder erst nach Zahlen? Gurau streut das so irgendwie rein, aber ich halte das für nicht sinnvoll.
\end{query}



\section{Gruppen und Körper}

\begin{definition}[Innere Verknüpfung]
    Sei $M \neq \e$ eine Menge.

    Eine \underline{innere Verknüpfung} auf $M$ ist eine Abbildung $\star: M \times M \arr M$.

    Wir schreiben $a \star b = \star(a, b)$.
\end{definition}

\begin{example}
    iwas triviales
    plus, mal
    Verkettung

    nicht-Beispiel: $\N$ mit minus
\end{example}

\begin{definition}[Halbgruppe]
    Ein Tupel $(M, \star)$ aus einer Menge $M$ und einer inneren Verknüpfung heißt \underline{Halbgruppe}, falls $\star$ assoziativ ist:

    $$\forall a, b, c \in M: (a \star b) \star c = a \star (b \star c).$$

    Dabei bedeuten die Klammern $()$, dass der Ausdruck innerhalb der Klammern zuerst ausgewertet wird. Wenn die Reihenfolge egal ist, dann lässt man die Klammern auch manchmal weg.
\end{definition}

\begin{example}
    $\N$ mit plus oder mal
    $\R$ mit plus oder mal
    Pot mit $\cap$
    
    was nicht assoziatives? --> vielleicht a hoch b (Matrixmult zu früh)
\end{example}

\begin{definition}[neutrales Element]
    Sei $(M, \star)$ eine Halbgruppe.

    $e \in M$ heißt \underline{neutrales Element}, falls gilt: $\forall a \in M: a \star e = a = e \star a$.

    Eine Halbgruppe, in der ein neutrales Element existiert, heißt \underline{Monoid}.
\end{definition}

\begin{lemma}[neutrales eindeutig]
    Sei $(M, \star)$ ein Monoid. Dann ist das neutrale Element eindeutig.
\end{lemma}
\begin{proof}
    Seien $e_1, e_2$ neutrale Elemente. Dann gilt:

    $e_1 = e_1 \star e_2 = e_2$.

    Dabei haben wir im ersten Schritt ausgenutzt, dass $e_2$ neutrales Element ist, und im zweiten Schritt, dass $e_1$ neutrales Element ist.
\end{proof}

\begin{example}
    $\R$ mit plus oder mal

    kein Monoid: $\N \backslash 0$ mit plus
\end{example}

\begin{definition}[Gruppe]
    Sei $(M, \star)$ ein Monoid mit neutralem Element $e$.
    \begin{enumerate}
        \item $a \in M$ heißt \underline{invertierbar}, falls $\exists b \in M: a \star b = e = b \star a$. $b$ heißt die \underline{Inverse} zu $a$.
        \item $(M, \star)$ heißt \underline{Gruppe}, falls jedes Element in $M$ invertierbar ist.
    \end{enumerate}
\end{definition}

\begin{lemma}[inverse eindeutig]
    Sei $(M, \star)$ ein Monoid mit neutralem Element $e$.

    Falls $a \in M$ invertierbar ist, dann ist die Inverse zu $a$ eindeutig.
\end{lemma}
\begin{proof}
    Seien $b_1, b_2$ Inverse zu $a$. Dann gilt:

    $b_1 = b_1 \star e = b_1 \star (a \star b_2) = (b_1 \star a) \star b_2 = e \star b_2 = b_2$.

    Dabei haben wir im zweiten Schritt verwendet, dass $b_2$ Inverse zu $a$ ist, im dritten Schritt die Assoziativität, und im vierten Schritt, dass $b_1$ Inverse zu $a$ ist.
\end{proof}

Im Allgemeinen schreibt man $b = a\m$. Man schreibt auch $b = -a$, falls $\star$ eine Addition ist, und $b = \frac{1}{a}$, falls $\star$ eine Multiplikation ist.

\begin{example}
    Bsp für Inverse in nur eine Richtung: Funktion?
    
    $\Z$ mit plus
    $\Q \backslash 0$

    keine Gruppe: $\Z$ mit mal
\end{example}

\begin{lemma}[Rechenregeln für Gruppen]
    Sei $(M, \star)$ eine Gruppe mit neutralem Element $e$.
    
    Dann gilt $\forall a, b, c \in M$:
    \begin{enumerate}
        \item \underline{Kürzen}: $a \star b = a \star c \implies b = c$, und analog $a \star c = b \star c \implies a = b$.
        \item $a \star b = e \implies b = a\m$.
        \item $(a\m)\m = a$.
        \item \underline{Regel von Hemd und Jacke}: $(a \star b)\m = b\m \star a\m$.

        Der Name kommt daher, dass man beim Anziehen zuerst das Hemd und dann die Jacke anzieht, aber beim Ausziehen zuerst die Jacke und dann das Hemd.
    \end{enumerate}
\end{lemma}
\begin{proof}
    Seien $a, b, c \in M$ beliebig.
    \begin{enumerate}
        \item Es gelte $a \star b = a \star c$. Da $M$ eine Gruppe ist, existiert $a\m$, und wir können es von links an die Gleichung anknüpfen:

        $a\m \star (a \star b) = a\m \star (a \star c)$. Mit Assoziativität folgt:

        $b = e \star b = (a\m \star a) \star b = (a\m \star a) \star c = e \star c = c$.

        Die andere Kürzungsregel folgt analog.
        \item Es gelte $a \star b = e$. Es gilt auch $e = a \star a\m$. Aus (i) folgt $b = a\m$.
        \item Per Definition gilt $a \star a\m = e = a\m \star a$. Nach derselben Definition, nur andersherum gelesen, bedeutet das aber auch, dass $a$ die Inverse zu $a\m$ ist, d.h. $a = (a\m)\m$.
        \item Es gilt: $(a \star b) \star (b\m \star a\m) = a \star (b \star b\m) \star a\m = a \star e \star a\m = a \star a\m = e$.

        Analog gilt auch: $(b\m \star a\m) \star (a \star b) = b\m \star e \star b = e$.

        Per Definition ist damit $(b\m \star a\m) = (a \star b)\m$.
    \end{enumerate}
\end{proof}

\begin{definition}[abelsche Gruppe]
    $(M, \star)$ Halbgruppe, Monoid, oder Gruppe, heißt \underline{kommutativ} oder \underline{abelsch}, falls gilt:

    $$\forall a, b \in M: a \star b = b \star a.$$
\end{definition}

Manchmal hilft eine dumme Eselsbrücke:
\begin{align*}
    & \text{\textbf{K}} \quad \text{ommutativ}\\
    & \text{\textbf{A}} \quad \text{ssoziativ}\\
    & \text{\textbf{\,I}} \quad \, \text{nverse}\\
    & \text{\textbf{N}} \quad \text{eutrales Element}\\[0.5\baselineskip]
    \implies & \text{\textbf{A B E L}}\\
\end{align*}

\begin{query}
    Darf ich "dumme Eselsbrücke" schreiben?

    Wir mein Oberstufen Physiklehrer immer gesagt hat: Je dümmer die Brücke, desto leichter kommt der Esel drüber...
\end{query}

\begin{example}
    $\R$ mit plus
    Permutationen

    keine abelsche Gruppe: Matrixmultiplikation?
    Funktionverknüpfung!
    $(\{f: \R \arr \R | f \text{ ist bijektiv}\}, \circ)$ ist eine nichtabelsche Gruppe:
    \begin{itemize}[label=]
        \item $(h \circ (g \circ f))(x) = (h(g(f(x)))) = ((h \circ g) \circ f)(x) \implies$ assoziativ
        \item Neutrales Element: $id_{\R}: x \mapsto x$
        \item Inverse $f\m$ existiert nach \ref{lem-existence-inverse-function}, da $f$ bijektiv ist
        \item Aber Verknüpfung von Funktionen ist i.A. nicht kommutativ, z.B. $f, g: \R \arr \R$ mit $f(x) = 2x$ und $g(x) = x + 1$ (beide sind bijektiv), dann ist $(g \circ f)(x) = 2x + 1$ aber $(f \circ g)(x) = 2x + 2$.
    \end{itemize}
\end{example}

\begin{query}
    Faktorgruppe? d.h. $\faktor{G}{N} = \{a \star N\}$ mit $N \subset G$

    Braucht:

    - Untergruppe

    - Normalteiler
    
    ist es das wert?
\end{query}
Bis jetzt haben wir uns Mengen mit einer Verknüpfung angeschaut. Viele relevante Mengen (in erster Linie Zahlen, aber später werden wir auch noch weitere Beispiele kennenlernen) haben allerdings zwei elementare Verknüpfungen. Das führt zur Definition eines Körpers (engl. \emph{field}), wofür unser motivierendes Beispiel die reellen Zahlen mit Addition und Multiplikation sind.

\begin{definition}[Ring]
    Sei $(R, +, \cdot)$ ein Tupel aus einer Menge $R$, und zwei inneren Verknüpfungen $+(a, b) = a + b$ ("Addition") und $\cdot(a, b) = a \cdot b$ ("Multiplikation").
    
    Das Tupel heißt \underline{Ring}, falls gilt:
    \begin{enumerate}[label=(R\arabic*)]
        \item $(R, +)$ ist kommutative Gruppe
        \item $(R, \cdot)$ ist Halbgruppe
        \item Es gelten die Distributivgesetze: $\forall a, b, c \in R$ ist
        \begin{itemize}[label=]
            \item $a \cdot (b + c) = a \cdot b + a \cdot c$
            \item $(a + b) \cdot c = a \cdot c + b \cdot c$
        \end{itemize}
        Um Klammern zu sparen, definieren wir hierbei, dass Multiplikation vor Addition ausgeführt wird.
    \end{enumerate}
    Ein Ring heißt \underline{kommutativ}, falls $(R, \cdot)$ kommutative Halbgruppe ist.

    Ein Ring heißt \underline{Ring mit Eins}, falls $(R, \cdot)$ ein Monoid ist.
\end{definition}

\begin{shadedbox}
    \textbf{Notation:}

    Das neutrale Element der Addition nennen wir $0$ oder $0_R$.
    
    Das neutrale Element der Multiplikation, falls es existiert, nennen wir $1$ oder $1_R$.
    
    Die Inverse zu $a \in R$ bezüglich der Addition bezeichnen wir mit $-a$, und wir schreiben $a + (-b) = a - b$.
    
    Die Inverse zu $a$ bezüglich der Multiplikation, falls diese existiert, bezeichnen wir diese mit $a\m$.
    
    Wir bemerken allerdings, dass die Schreibweise $a \cdot b\m = \frac{a}{b} = b\m \cdot a$ nur in einem kommutativen Ring Sinn ergibt.
\end{shadedbox}

\begin{example}
    Nullring hahaha
    
    $\Z$, $\Q$, $\R$ sind komm R mit Eins

    $2\Z$ Ring ohne Eins

    $f, +, \circ$ nicht komm Ring mit Eins? wenn $f$ Gruppem Homo
    $g \circ (f_1 + f_2)(x) = g(f_1(x) + f_2(x))$
    zu viel Einführungsarbeit?
\end{example}

\begin{lemma}[lemma]
    Sei $(R, +, \cdot)$ ein Ring. Dann gelten $\forall a, b \in R$ folgende Rechenregeln:
    \begin{enumerate}
        \item $0_R \cdot a = 0_R = a \cdot 0_R$
        \item $a \cdot (-b) = - (a \cdot b) = (-a) \cdot b$
        \item Falls $R \neq \{0_R\}$ Ring mit Eins ist, dann ist $0_R \neq 1_R$.
    \end{enumerate}
\end{lemma}

Zwei äquivalente Wege:

\begin{definition}[Körper]
    Sei $(K, +, \cdot)$ ein Tupel aus einer Menge $K$, und zwei inneren Verknüpfungen $+(a, b) = a + b$ ("Addition") und $\cdot(a, b) = a \cdot b$ ("Multiplikation").
    
    Das Tupel heißt \underline{Körper}, falls gilt:
    \begin{enumerate}[label=(K\arabic*)]
        \item $(K, +)$ ist kommutative Gruppe mit neutralem Element $0_K$
        \item $(K \backslash \{0_K\}, \cdot)$ ist kommutative Gruppe mit neutralem Element $1_K$
        \item Es gelten die Distributivgesetze: $\forall a, b, c \in R$ ist
        \begin{itemize}[label=]
            \item $a \cdot (b + c) = a \cdot b + a \cdot c$
            \item $(a + b) \cdot c = a \cdot c + b \cdot c$
        \end{itemize}
    \end{enumerate}
    In der Literatur findet man oft die Schreibweise $K^* = K \backslash \{0_K\}$.
\end{definition}

\begin{definition}[Körper]
    Ein Ring $(K, +, \cdot)$ heißt \underline{Körper}, falls zusätzlich gilt:
    \begin{enumerate}[label=(K\arabic*)]
        \item $(K, +, \cdot)$ ist kommutativer Ring mit Einselement $1_K$ und Nullelement $0_K$
        \item $1_K \neq 0_K$
        \item $\forall a \in K \backslash \{0_K\}$ existiert eine multiplikative Inverse $a\m \in K$
    \end{enumerate}
    In der Literatur findet man oft die Schreibweise $K^* = K \backslash \{0_K\}$.
\end{definition}

\begin{example}
    - Aus (K2) folgt bereits, dass $0_K \neq 1_K$, d.h. ein Körper muss mindestens zwei Elemente haben. Dementsprechend kann der Nullring kein Körper sein.

    - $\Q$, $\R$ Körper, $\Z$ kein Körper

    - Restklassenring i.A. kein Körper: z.B. in $\Z_4$ hat $[2]$ keine mult Inverse
\end{example}




\section{Zahlen}
\label{sec-numbers}

Zahlen sind intrinsisch motiviert. Zuerst einmal sind die natürlich Zahlen $\N$ ein mathematisches Mittel, um Objekte in der physischen Welt zu zählen\footnote{Man kann auch nur Objekte in der mathematischen Welt zählen, wenn man so sehr darin gefangen steckt.}. Diese werden mit negativen Zahlen zu den ganzen Zahlen $\Z$ erweitert, um Prozesse der Subtraktion zu beschreiben. Hier haben wir Subtraktion als inverse Operation zur Addition eingeführt, wodurch negative Zahlen mathematisch von der Forderung der Abgeschlossenheit bezüglich Subtraktion kommen, sie haben aber auch eine nichtmathematische Motivation (z.B. Schulden). Rationale Zahlen $\Q$ sind nichtmathematisch die Konsequenz davon, Zahlen für physische Messungen zu verwenden, für die ganze Zahlen nicht präzise genug sind, oder mathematisch aus der Forderung der Existenz von multiplikativen Inversen.

Relle Zahlen $\R$ haben wenig nichtmathematische Motivation. Auch wenn die Existenz von nicht rationalen (irrationalen) Zahlen wie $\sqrt{2}$ oder $\pi$ bereits seit der Antike bekannt ist, haben diese ihre Relevanz doch eher als mathematischen Vervollständigung von $\Q$, als als Messergebnisse in der physischen Welt. Die komplexen Zahlen $\C$ sind rein mathematisch motiviert als algebraische Vervollständigung von $\R$.

\begin{definition}[Natürliche Zahlen]
    \label{def-natural-numbers}
    Intuitiv kann man die natürlichen Zahlen fordern, indem man $0$ als "nichts" definiert, $1$ als "eines", und von da an jede kommende Zahl als "die Zahl davor $+ 1$". Mathematisch wäre streng genommen nichts falsch damit. Der einzige ästhetische Mangel ist, dass man dazu ein zusätzliches Axiom bräuchte, was aber nicht nötig ist, da man die intrinsisch bekannte Struktur der natürlichen Zahlen mit den bereits bekannten Axiomen konstruieren kann. Das macht man genauso, wie beschrieben, d.h. man fängt mit "nichts" an und geht jeweils Schritt für Schritt weiter "hoch". Dafür verwendet man die Existenz der leeren Menge (ZF5) sowie ein bisschen Trixerei mit Mengenklammern. Dann definiert man sich die \underline{natürlichen Zahlen} über eine Identifikation:
    \begin{itemize}[label=]
        \item $0 = \e$ für $n = 0$,
        \item der \underline{Nachfolger} von $n$ ist definert als $n \cup \{n\}$,
        \item $\N$ ist definiert als die kleinste Menge, die $0$ und alle Nachfolger enthält.
    \end{itemize}
    Dann ist z.B. $1 = \e \cup \{\e\} = \{\e\}$, $2 = \{ \e, \{\e\} \}$, $3 = \{ \e, \{ e \}, \{\e, \{\e\}\} \}$ etc. Dann ist $n$ definiert als die Menge, die genau $n$ Elemente enthält.
    
    Wir bemerken, dass $\N$ nach (ZF7) existiert und wohldefiniert ist.
\end{definition}

\begin{definition}[Ganze Zahlen]
    $\Z$
\end{definition}

\begin{definition}[Rationale Zahlen]
    $\Q$
\end{definition}

\begin{proposition}[$\Q$ ist abzählbar]
    
\end{proposition}

\begin{definition}[Reelle Zahlen]
    $\R$
\end{definition}

\begin{proposition}[Existenz von Suprema in $\R$]
    
\end{proposition}

\begin{definition}[Kardinalität]
    \label{def-cardinality}
    Sei $A$ eine Menge. Dann bezeichnen wir mit $|A|$ die \underline{Kardinalität} (oder \underline{Mächtigkeit}) von $A$. Wir nennen $A$:
    \begin{itemize}
        \item \underline{endlich} mit $|A| = n$, falls $A \cong \{1, ..., n\}$ mit $n \in \N$, oder mit $|A| = 0$ falls $A = \e$.
        \item \underline{abzählbar unendlich}, falls $A \cong \N$.
        \item \underline{überabzählbar}, sonst.
    \end{itemize}
\end{definition}

\begin{theorem}[$\Q$ ist abzählbar]
    \label{theorem-cantor-1}

\end{theorem}

\begin{theorem}[$\R$ ist überabzählbar]
    \label{theorem-cantor-2}

\end{theorem}

\begin{proposition}[$\R$ erfüllt die Archimedische Eigenschaft]
    
\end{proposition}

\begin{theorem}[$\Q$ liegt dicht in $\R$]
    
\end{theorem}

\begin{definition}[Unendlich]
    $$|A| = \infty bEcAuSeInFiNiTyIsAnUmBeR!1!$$
\end{definition}

\begin{proposition}[Satz von Cantor]
    don't go insane!
\end{proposition}

\begin{definition}[Komplexe Zahlen]
    $\C$
\end{definition}

\begin{definition}[Eigenschaften von komplexen Zahlen]
    \begin{enumerate}
        \item re, im
        \item cc
        \item Betrag
        \item Winkel, Radialdarstellung
    \end{enumerate}
\end{definition}

\begin{lemma}[lemma]
    mult inverses
\end{lemma}

\begin{definition}[Potenz, Wurzel]
    
\end{definition}



\section{GURAU}

\begin{definition}[ordinals, cardinals]
    vlt soll ich einfach auf das Originalscript verweisen?
\end{definition}

\chapter{Topologie}

\begin{query}
    Wie wäre es, statt am Anfang jedes Kapitels ein Inhaltsverzeichnis einzufügen, dass ich ein Stichwortverzeichnis der wichtigsten Definitionen, Lemmata, und Sätze gebe?

    Also Bsp mal angenommen das Kapitel wäre (Auszug aus Mengen und Zahlen, 1 von 6 Unterkapiteln):
    \begin{itemize}
        \item Innere Verknüpfung
        \item Halbgruppe
        \item Monoid
        \item Lemma: neutrales Element eindeutig
        \item Gruppe
        \item Lemma: Inverse eindeutig
        \item Lemma: Rechenregeln für Gruppen
        \item Abelsche Gruppe
        \item Gruppenhomomorphismus
        \item Lemma: Eigenschaften von Gruppenhomomorphismen
        \item Ring
        \item Lemma: Rechenregeln in Ringen
        \item Körper
    \end{itemize}
    Dann sähe meine Übersicht vlt so aus (jeweils mit Links):
    \begin{itemize}
        \item Gruppe
        \item Lemma: Rechenregeln für Gruppen
        \item Abelsche Gruppe
        \item Gruppenhomomorphismus
        \item Lemma: Eigenschaften von Gruppenhomomorphismen
        \item Körper
    \end{itemize}

    Dann müsste ich mich aber entscheiden, was wichtig genug ist um auf das Verzeichnis zu kommen und was nicht, und das traue ich mir nicht zu. Zum Beispiel bin ich jetzt schon wenn ich mit dem Beispiel fertig bin überzeugt dass das keine gute Idee ist. Und das war noch eins der "klarsten", im Kapitel über Zahlen wird es viel schwerer Sachen rauszukürzen.

    Aber ich lasse die Frage und das Beispiel trotzdem mal drin, weil es einen guten Eindruck davon vermittelt, wie lang so ein Inhaltsverzeichnis am Kapitelanfang wäre.
\end{query}

some motivational words

\section{Grundlegende Begriffe}

Längere Motivation: Offene Intervalle

\begin{definition}[Topologie]
    Sei $X$ eine Menge, und $\T \subset \Pot(X)$ ein Mengensystem, d.h. eine Menge, deren Elemente Teilmengen von $X$ sind. Wir können also Teilmengen $U \subset X$ danach bewerten, ob sie Elemente in $\T$ oder nicht sind.
    
    $\T$ heißt \underline{Topologie} und das Tupel $(X, \T)$ \underline{topologischer Raum}, falls gilt:
    \begin{enumerate}[label=(T\arabic*)]
        \item $\e, X \in \T$
        \item $(U_{\alpha})_{\alpha \in I} \subset \T \implies \bigcup\limits_{\alpha \in I} U_{\alpha} \in \T$
        \item $(U_i)_{i \in \llbracket 1, n \rrbracket} \subset \T \implies \bigcap\limits_{1 = 1}\limits^n U_i \in \T$
    \end{enumerate}
    Wie die Indizes andeuten, kann dabei die Vereinigung in (T2) beliebig (auch überabzählbar) sein, der Schnitt in (T3) aber höchstens endlich.

    Mengen $U \in \T$ nennen wir \underline{offen}.
\end{definition}

\begin{example}
    \begin{enumerate}
        \item Die Menge aller offenen Intervalle auf $\R$ bildet eine Topologie. Dabei zählen wir Mengen der Form $(a, \infty)$, $(-\infty, b)$, und $(- \infty, \infty) = \R$ auch als offene Intervalle.
        \item Für jede Menge $X$ kann man zwei triviale Topologien definieren:
        \begin{itemize}
            \item Die \underline{diskrete Topologie} $\T_{disk} = \Pot(X)$, d.h. jede Menge ist offen.
            \item Die \underline{Klumpentopologie} $\T_{triv} = \{\e, X\}$ (\underline{triviale Topologie}, \underline{indiskrete Topologie}), d.h. nur die leere Menge und die "ganze Menge" sind offen.
        \end{itemize}
        Beide erfüllen trivialerweise (T1-T3).
        \item $\T = (\e, \{0\}, \R^*, \R)$ ist eine Topologie auf $\R$:
        \begin{itemize}%[label=]
            \item $\e \cup \{0\} = \{0\} \in \T$, $\{0\} \cup \R^* = \R \in \T$, usw
            \item $\R \cap \{0\} = \{0\} \in \T$, $\R^* \cap \{0\} = \e \in \T$, usw
        \end{itemize}
        Die letzten drei Beispiele waren zugegebenermaßen eher konstruiert und weniger anwendungsbezogen. Das nächste Beispiel ist allerdings sehr relevant:
        \item Die Menge aller offenen Intervalle $(a, b)$ in $\R$ bildet eine Topologie. Dabei zählen wir Mengen der Form $(a, \infty)$, $(-\infty, b)$, und $(- \infty, \infty) = \R$, sowie disjunkte Vereinigungen von offenen Intervallen, auch als offene Intervalle.
        
        Der Leser kann sich (z.B. durch Skizzen) selbst überzeugen, dass es sich um eine Topologie handelt. Formal werden wir (T1-T3) hier nicht nachprüfen, aber wir nehmen schonmal vorweg, dass es sich um einen Spezialfall von metrischer Topologie handelt. Damit folgt die Behauptung aus \ref{lem-metric-induces-topology}.
    \end{enumerate}
\end{example}

MEHR BEISPIELE?

\begin{definition}[offene Umgebung]
    Sei $(X, \T)$ ein topologischer Raum und $x \in X$.

    Eine \underline{offene Umgebung} von $x$ ist eine Menge $U \in \T$ mit $x \in U$.
\end{definition}

ein paar Worte über Bedeutung auf $\R^n$: immer größer als $x$, aber trotzdem beliebig klein.

Feinere / gröbere Topologien?

Basis hier? oder unten?

\begin{definition}[Hausdorff Eigenschaft]
    Ein metrischer Raum $(X, \T)$ heißt \underline{Hausdorff-Raum}, falls man zwei beliebige Punkte durch offene Mengen trennen kann:

    $\forall x, y \in X$ mit $x \neq y$ $\exists$ offene Umgebungen $U_x, U_y \in T$ von $x, y$ mit $U_x \cap V_y = \e$.
\end{definition}

\begin{example}
    $\T_{disk}$ HD, $\T_{triv}$ nicht HD

    offene Intervalle HD

    Das Problem von Räumen, die nicht Hausdorff sind:
    Zwei Kopien von $\R$...
\end{example}

\begin{definition}[abgeschlossene Mengen]
    Sei $(X, \T)$ ein topologischer Raum. Eine Menge $A \subset X$ heißt \underline{abgeschlossen}, falls $A^C$ offen ist.
\end{definition}

\begin{definition}[Unterraum]
    Sei $(X, \T)$ ein topologischer Raum, und $A \subset X$.

    Dann heißt $\T_A = \{U \cap A | U \in \T\}$ die \underline{Unterraumtopologie}, die von $X$ auf $A$ induziert wird.

    Mengen $U \in \T_A$ nennen wir \underline{offen in $A$} (oder \underline{relativ offen}).
    
    Mengen $A$ mit $A^C \in \T_A$ nennen wir \underline{abgeschlossen in $A$} (oder \underline{relativ abgeschlossen}).
\end{definition}

\begin{exercise}
    Zeigen Sie, dass $\T_A$ wieder eine Topologie ist.
\end{exercise}

\begin{definition}[accumulation point, innerer Punkt, isolated point]
    $\equiv$

    test$\iff \hspace{-17pt} \not \hspace{17pt}$test
\end{definition}

\begin{definition}[Inneres, Abschluss, Rand]
    Sei $(X, \T)$ ein topologischer Raum, und $U \subset X$.
\end{definition}

- Beispiele mit trivialen Topologien

\begin{lemma}[Inneres ist offen]
    
\end{lemma}

\begin{lemma}[Abschluss ist abgeschlossen]
    
\end{lemma}

\begin{corollary}[Rand]
    trivial
\end{corollary}

Warnungen? "we stress that..."

\begin{definition}[zusammenhängend]
    Sei $(X, \T)$ ein topologischer Raum.
    
    Eine Menge $Y \subset X$ heißt \underline{zusammenhängend}, falls man sie nicht in offene disjunkte Mengen teilen kann:

    $\not\exists U, V \in \T: Y \subset U \cup V$ und $U \cap V = \e$.
\end{definition}

\begin{definition}[wegzusammenhängend]
    Vielleicht erst später? vlt in einem Kapitel über Kurven in $\R^n$?
\end{definition}

\begin{lemma}[wegzusammenhängend impliziert zusammenhängend]
    
\end{lemma}

\begin{definition}[kompakt]
    Sei $(X, \T)$ ein topologischer Raum.

    Eine Menge $C \subset X$ heißt \underline{kompakt}, falls für jede offene Überdeckung von $C$ eine endliche Teilüberdeckung existiert:

    $\forall (U_{\alpha})_{\alpha \in I}$ mit $C \subset \bigcup\limits_{\alpha \in I} U_{\alpha}$ $\exists \alpha_1, ..., \alpha_n \in I: C \subset \bigcup\limits_{i=1}^{n} U_{\alpha_i}$.
\end{definition}

\begin{lemma}[abgeschlossene Teilmengen kompakter Mengen sind kompakt]
    zwei Aussagen
\end{lemma}

\begin{proposition}[Punktkompaktifizierung]
    Sei $(X, \T)$ ein topologischer Raum, und $Y \subset X$ beliebig.

    Wir können aus $Y$ durch Hinzunahme eines einzigen Punktes, den wir $\infty$ nennen\footnote{Dabei soll/muss $\infty$ nicht das bekannte "unendlich" sein, sondern ein beliebiger Punkt, der nicht in $X$ liegt. Dass man diesen auch mit $\infty$ bezeichnet, kommt daher, dass er bezüglich der neu kontruierten Topologie von $X \cup \{\infty\}$ als "unendlich weit weg" von $X$ interpretiert werden kann.}, eine kompakte Menge $\bar{Y} = Y \cup \{\infty\}$ machen.
\end{proposition}




\section{Folgentopologie und Metrische Topologie}

\textbf{Folgentopologie}

\begin{query}
    Hierfür brauche ich mindestens Konvergenz von Folgen. Wo soll ich das einführen? Optionen die ich sehe:
    \begin{enumerate}[label=\thequery.\arabic*.]
        \item Find ich am sinnvollsten: Kapitel Topologie und Folgen zu einem Riesenkapitel mergen und Stoff relativ zueinander umordnen dass es halbwegs aufeinander aufbaut
        \item Finde ich am sinnvollsten 2: Kapitel Topologie und Folgen vertauschen
        \item Konvergenzbegriff hier einführen, kurz erklären mit Beispielen, dann mit Folgentopologie weitermachen \& den Rest von Folgen im nächsten Kapitel machen
        \item In Kapitel "Mengen und Zahlen" werden sowieso Familien und Folgen definiert (aber nichts größer damit gemacht), da könnte ich auch eine Definition von Konvergenz von Folgen dazu packen
        \item Halte ich für nicht sinnvoll: Folgentopologie hier rauslassen und in Kapitel Folgen nachholen
    \end{enumerate}
\end{query}

\begin{definition}[Folgentopologie]
    abgeschlossen bzgl eines Konvergenzbegriffes
\end{definition}

\begin{definition}[Folgenkompakt]
    
\end{definition}

noch Sachen aus Gurau Kapitel 4?



\textbf{Metrische Topologie}


\begin{definition}[Metrischer Raum]
    
\end{definition}

\begin{example}
    \begin{enumerate}
        \item Dirac Metrik
        \item Betrag auf $\R$
        \item Betrag auf $\R^n$
        \item $p$-Norm und $\infty$-Norm auf $\R^n$
    \end{enumerate}
\end{example}

\begin{definition}[offener Ball]
    Sei $(X, d)$ ein metrischer Raum, und $x \in X$.

    Der \underline{offene Ball} mit Radius $r$ um den Punkt $x$ ist die Menge $B_r(x) = \{y \in X | d(x,y) < r\}$.

    Der \underline{abgeschlossene Ball} mit Radius $r$ um den Punkt $x$ ist die Menge $\bar{B}_r(x) = \{y \in X | d(x,y) \leq r\}$.
\end{definition}

\begin{proposition}[Metrik induziert Topologie]
    \label{lem-metric-induces-topology}
    Sei $(X, d)$ ein metrischer Raum.

    Dann trägt $X$ eine kanonisch induzierte Topologie $\T_d$ definiert durch $U \in \T_d \iff \forall x \in U \exists \epsilon > 0: B_{\epsilon}(x) \subset U$.
\end{proposition}

\begin{lemma}[metrischer Rand]
    Sei $(X, d)$ ein metrischer Raum mit induzierter Topologie $\T_d$.

    Dann gilt bezüglich dieser Topologie
    \begin{enumerate}
        \item Für $Y \subset X$: $y \in \partial Y \iff \forall \epsilon > 0$
        \item Für Bälle: $B_r(x)$ ist offen ($\implies B_r(x)^o = B_r(x)$)
        \item Für Bälle: $\bar{B}_r(x)$ ist abgeschlossen ($\implies \bar{B}_r(x) = \bar{B_r(x)}$)
        \item muss schauen ob ich noch mehr von diesen $\partial Y = \bar{Y} \backslash Y^o$ etc Miniaussagen reinpacken will, und wie viele davon spezifisch für metrische Räume sind
    \end{enumerate}
\end{lemma}
\begin{proof}
    Hausaufgabe?
\end{proof}

\begin{example}
    $\R^2$ mit Halbebene $\{x > 0\}$ ist ein MEGA gutes Beispiel
\end{example}

\begin{lemma}[Metrisch $\implies$ Hausdorff]
    Sei $(X, d)$ ein metrischer Raum mit induzierter Topologie $\T_d$.

    Dann ist $(X, \T_d)$ ein Hausdorr-Raum.
\end{lemma}
\begin{proof}
    Seien $x, y \in X$, $x \neq y$.

    Wir müssen zwei offene, nichtleere Mengen finden, die $x$ und $y$ trennen. Dazu verwenden wir die Metrik:

    Wir setzen $d = d(x,y)$, dann ist $d > 0$ wegen (M1).

    Dazu betrachten wir die beiden offenen Bälle $B_{\frac{d}{2}}(x)$ und $B_{\frac{d}{2}}(y)$.

    Offensichtlich gilt $B_{\frac{d}{2}}(x) \neq \e \neq B_{\frac{d}{2}}(y)$, da $d(x, x) = 0 < d \implies x \in B_{\frac{d}{2}}(x)$ und analog für $B_{\frac{d}{2}}(y)$.

    Jetzt wollen wir zeigen, dass der Schnitt leer ist. Dazu nehmen wir an, es gäbe ein $z \in B_{\frac{d}{2}}(x) \cap B_{\frac{d}{2}}(y)$. Dieses $z$ kann es offensichlicht nicht geben, da es näher als $\frac{d}{2}$ an $x$ und näher als $\frac{d}{2}$ an $y$ liegen müsste.
    
    Das formalisieren wir mit der Dreiecksungleichung: $d = d(x, y) \leq d(x, z) + d(z, y) < \frac{d}{2} + \frac{d}{2} = d$.\absurdum

    Also gibt es $z$ nicht und $B_{\frac{d}{2}}(x) \cap B_{\frac{d}{2}}(y) = \e$.

    Damit haben wir unsere Mengen gefunden.
\end{proof}

\begin{query}
    nicht sicher ob das hier hingehört, aber Gurau macht das
\end{query}

\begin{definition}[Norm]
    \begin{enumerate}[label=(N\arabic*)]
        \item positiv definit
        \item Skalare mit Betrag rausziehen
        \item Dreiecks ungl
    \end{enumerate}
\end{definition}

\begin{definition}[Normäquivalenz]
    Zwei Normen $|\cdot|_1$, $|\cdot|_2$ heißen \underline{äquivalent}, falls $\exists c, C > 0: c |\cdot|_1 \leq |\cdot|_2 \leq |C \cdot|_1$.
\end{definition}

\begin{example}
    
\end{example}

\begin{proposition}[Auf $\R^n$ sind alle Normen äquivalent]
    Es reich z.z. dass alle Normen äquivalent zur Standardnorm sind.
\end{proposition}

Basis hier? oder oben?

\begin{lemma}[$\Q$ liegt dicht in $\R$]
    Bezüglich der Standardtopologie gilt $\bar{\Q} = \R$.
\end{lemma}



\section{Stetigkeit}

\begin{definition}[$\epsilon$-$\delta$-Stetigkeit]
    $$\forall \epsilon > 0 \exists \delta > 0 \forall x \in D: |x-x_0| < \delta \implies |f(x)-f(x_0)| < \epsilon$$
\end{definition}

- Erklärung
- Zwei Bilder

\begin{definition}[topologische Stetigkeit]
    Seien $(X, \T_X)$ und $(Y, \T_Y)$ topologische Räume, und $f: X \arr Y$ eine Funktion.

    $f$ heißt stetig, falls gilt: $V \in \T_Y \implies f\m(V) \in \T_X$.
\end{definition}

\begin{definition}[Folgenstetigkeit]
    $f: X \arr Y$ heißt \underline{folgenstetig}, falls $\forall (x_n)_{n \in \N}$ mit Grenzwert in $X$ gilt: $\lim_{n \arr \infty} f(x_n) = f(\lim_{n \arr \infty}x_n)$.
\end{definition}

\begin{lemma}[stetige Bilder kompakter Mengen sind kompakt]
    
\end{lemma}

\begin{lemma}[Metrik ist stetig]
    
\end{lemma}

\begin{remark}
    flow chart ansprechen, Sorgenfrey Topologie
\end{remark}

\begin{definition}[Gleichmäßige Stetigkeit]
    $$\forall \epsilon > 0 \exists \delta > 0 \forall x, y \in D: |x-y| < \delta \implies |f(x)-f(y)| < \epsilon$$
\end{definition}

\begin{lemma}[gleichmäßig stetig $\implies$ stetig]
        
\end{lemma}

\begin{example}
    \begin{enumerate}
        \item $id$ ist gleichmäßig stetig
        \item $f: \R_+ \arr \R_+, \quad f(x) = \frac{1}{x}$ ist stetig aber nicht gleichmäßig stetig
    \end{enumerate}
\end{example}

\begin{definition}[Lipschitz Stetigkeit]
    $$\exists L \leq 0: \forall x, y \in D: |f(x)-f(y)| < L \cdot |x-y|$$
\end{definition}

\begin{lemma}[Lipschitz stetig $\implies$ gleichmäßig stetig]
    
\end{lemma}

\begin{example}
    \begin{enumerate}
        \item $id$ ist Lipschitz stetig
        \item $f: \R_+ \arr \R_+, \quad f(x) = \sqrt{x}$ ist gleichmäßig aber nicht Lipschitz ($x=0$)
        \item Polynome sind Lipschitz auf $[a, b]$, nicht Lipschitz auf $\R$ (Problem $\infty$)
    \end{enumerate}
\end{example}

\chapter{auxies.aux}

Auxies!

\begin{definition}[Kaffee \hfill \textit{\small{vgl Gurau 9999999}}]
    gsrngrndguisrnfojesgsrgrnsuivnfdoheihrwig
\end{definition}

\newpage



\appendix


\backmatter
\printbibliography


\end{document}